\chapter{Complexity — Simulation}
%%% generated by the blueprint pipeline from TCSlib.Complexity.TuringMachine.Simulation
%%%%%%%%%%%%%%%%%%%%%%%%%%%%%%%%%%%%%%%%%%%%%%%%%%%%%%%%%%%%%%%%%%%%%%%%%%%%%%%
\section{Overview}

This module provides generic building blocks from which larger machine
constructions are assembled: fixed-word emission chains that write a word to
the output one symbol per step and then halt; control actions that only move
the input head and change state, together with a rewind-to-start scan of the
input head; disjoint tape-block extensions that run a machine on the left or
right block of a larger machine in exact lockstep; the branch union of two
machines selected by a Boolean; a normalization identity for work-tape writes;
and a buffered sequential simulator that runs a first machine into a middle
buffer tape, rewinds that buffer, and then serves the buffered word as the
virtual input of a second machine, with boundary clamping realized by a
Boolean arrival tag.

%%%%%%%%%%%%%%%%%%%%%%%%%%%%%%%%%%%%%%%%%%%%%%%%%%%%%%%%%%%%%%%%%%%%%%%%%%%%%%%
\section{Declarations}

\begin{lemma}[Work tapes after applying an action]
\lean{Turing.Action.apply_workTapes}
\label{Turing.Action.apply_workTapes}
\leanok
\uses{Turing.Action, Turing.Action.apply, Turing.Cfg, Turing.Cfg.workTapeSymbols}
\difficulty{2}
Let $a$ be a one-step action of a multi-tape Turing machine with $k$ work
tapes, let $c$ be a configuration of such a machine, and let $i$ be one of the
$k$ work-tape indices.  After applying $a$ to $c$, the content of the $i$-th
work tape equals the old content updated at the cell under the $i$-th work
head, where the new cell value is the symbol that $a$ proposes to write there
and, when $a$ declines to write on that tape, the symbol previously scanned:
an explicit blank write erases the cell, while no write at all copies the
scanned symbol back unchanged.  No finiteness or liveness hypothesis is
required.
\end{lemma}

\begin{definition}[One step of an emission chain]
\lean{Turing.FinTM.emitAction}
\label{Turing.FinTM.emitAction}
\leanok
\uses{Turing.Action}
Given a Boolean word $w$, a map $e$ from $\{0,\dots,\abs{w}\}$ into the states
of a machine with $k$ work tapes over the Boolean alphabet, and an index
$i \in \{0,\dots,\abs{w}\}$, the one-step action that, when $i < \abs{w}$,
appends the $i$-th symbol of $w$ (zero-indexed) to the output and passes to
state $e(i+1)$, and otherwise emits nothing and halts.  In both cases the
input head and all work tapes are left untouched.
\end{definition}

\begin{lemma}[Partial run of an emission chain]
\lean{Turing.FinTM.emit_run}
\label{Turing.FinTM.emit_run}
\leanok
\uses{Turing.Action.apply, Turing.Cfg, Turing.FinTM.emitAction, Turing.MultiTapeTM, Turing.MultiTapeTM.runFrom, Turing.MultiTapeTM.runFrom_succ_eq_step', Turing.MultiTapeTM.step}
\difficulty{4}
Let $M$ be a multi-tape Turing machine with $k$ work tapes over the Boolean
alphabet running on some input $x$, let $w$ be a Boolean word, and let $e$ be
a map from $\{0,\dots,\abs{w}\}$ into the states of $M$ such that from each
state $e(i)$, regardless of the symbols read, the transition of $M$ is the
emission action for $w$ at index $i$: while $i < \abs{w}$ it appends the
$i$-th symbol of $w$ to the output and passes to $e(i+1)$, and at
$i = \abs{w}$ it halts, in both cases leaving the input head and all work
tapes fixed.  If a configuration $c$ is in state $e(0)$, then for every
$t \le \abs{w}$, after $t$ steps from $c$ the machine is in state $e(t)$ and
its output is the output of $c$ followed by the first $t$ symbols of $w$.
\end{lemma}

\begin{lemma}[Emission chain halts after the word]
\lean{Turing.FinTM.emit_halts}
\label{Turing.FinTM.emit_halts}
\leanok
\uses{Turing.Action.apply, Turing.Cfg, Turing.FinTM.emitAction, Turing.FinTM.emit_run, Turing.MultiTapeTM, Turing.MultiTapeTM.runFrom, Turing.MultiTapeTM.runFrom_succ_eq_step', Turing.MultiTapeTM.step}
\difficulty{3}
Let $M$ be a multi-tape Turing machine with $k$ work tapes over the Boolean
alphabet running on some input $x$, let $w$ be a Boolean word, and let $e$ be
a map from $\{0,\dots,\abs{w}\}$ into the states of $M$ such that from each
state $e(i)$, regardless of the symbols read, the transition of $M$ is the
emission action for $w$ at index $i$ (append the $i$-th symbol of $w$ and
pass to $e(i+1)$ while $i < \abs{w}$, halt at $i = \abs{w}$, never moving the
input head or touching a work tape).  If a configuration $c$ is in state
$e(0)$, then after exactly $\abs{w} + 1$ steps from $c$ the machine has
halted and its output is the output of $c$ followed by all of $w$.
\end{lemma}

\begin{definition}[Input-head control action]
\lean{Turing.FinTM.controlAction}
\label{Turing.FinTM.controlAction}
\leanok
\uses{Turing.Action}
For a movement $m \in \{-1, 0, +1\}$ and an optional successor state $q$, the
one-step action of a machine with $k$ work tapes over the Boolean alphabet
that moves the input head by $m$ (subject to clamping at the two boundary
blanks), writes nothing on any work tape, moves no work head, emits no
output, and passes to $q$, halting when $q$ is absent.
\end{definition}

\begin{lemma}[Input symbol at a given position]
\lean{Turing.FinTM.inputSymbol_at}
\label{Turing.FinTM.inputSymbol_at}
\leanok
\uses{Turing.Cfg, Turing.Cfg.inputSymbol, Turing.inputSymbolInner}
\difficulty{4}
Let $c$ be a configuration of a multi-tape Turing machine over the Boolean
alphabet on input $x$, and let $i \le \abs{x}$ be a natural number such that
the input head of $c$ sits at position $i + 1$, in the convention where
position $0$ is the left boundary blank and position $j$ with
$1 \le j \le \abs{x}$ scans the $j$-th input symbol.  Then the symbol scanned
on the input tape is the optional $i$-th entry of $x$ (zero-indexed): the
symbol $x_i$ when $i < \abs{x}$, and the right boundary blank when
$i = \abs{x}$.
\end{lemma}

\begin{definition}[Left-block extension of an action]
\lean{Turing.FinTM.leftAction}
\label{Turing.FinTM.leftAction}
\leanok
\uses{Turing.Action}
Given natural numbers $k$ and $l$, a state renaming $f : S \to S'$, and a
one-step action $a$ of a machine with $k$ work tapes over the Boolean
alphabet, the action of a machine with $k + l$ work tapes that performs $a$
on the first $k$ work tapes, writes nothing and moves no head on the last $l$
tapes, keeps the input-head movement and the emitted output of $a$, and
renames the successor state along $f$, preserving halting.
\end{definition}

\begin{definition}[Right-block extension of an action]
\lean{Turing.FinTM.rightAction}
\label{Turing.FinTM.rightAction}
\leanok
\uses{Turing.Action}
Given natural numbers $k$ and $l$, a state renaming $f : S \to S'$, and a
one-step action $a$ of a machine with $l$ work tapes over the Boolean
alphabet, the action of a machine with $k + l$ work tapes that performs $a$
on the last $l$ work tapes, writes nothing and moves no head on the first $k$
tapes, keeps the input-head movement and the emitted output of $a$, and
renames the successor state along $f$, preserving halting.
\end{definition}

\begin{definition}[Left-block embedding of a configuration]
\lean{Turing.FinTM.leftCfg}
\label{Turing.FinTM.leftCfg}
\leanok
\uses{Turing.Cfg}
Given a state renaming $f : S \to S'$, a configuration $c$ of a machine with
$k$ work tapes over the Boolean alphabet on input $x$, and arbitrary contents
and head positions for $l$ further tapes, the configuration of a machine with
$k + l$ work tapes on the same input whose state is that of $c$ renamed along
$f$ (halted exactly when $c$ is halted), whose input head and output agree
with those of $c$, whose first $k$ work tapes and heads are those of $c$, and
whose last $l$ work tapes carry the given inactive contents and head
positions.
\end{definition}

\begin{definition}[Right-block embedding of a configuration]
\lean{Turing.FinTM.rightCfg}
\label{Turing.FinTM.rightCfg}
\leanok
\uses{Turing.Cfg}
Given a state renaming $f : S \to S'$, a configuration $c$ of a machine with
$l$ work tapes over the Boolean alphabet on input $x$, and arbitrary contents
and head positions for $k$ further tapes, the configuration of a machine with
$k + l$ work tapes on the same input whose state is that of $c$ renamed along
$f$ (halted exactly when $c$ is halted), whose input head and output agree
with those of $c$, whose last $l$ work tapes and heads are those of $c$, and
whose first $k$ work tapes carry the given inactive contents and head
positions, for instance data left behind by an earlier phase.
\end{definition}

\begin{lemma}[Action application commutes with the left embedding]
\lean{Turing.FinTM.leftCfg_apply}
\label{Turing.FinTM.leftCfg_apply}
\leanok
\uses{Turing.Action, Turing.Action.apply, Turing.Cfg, Turing.FinTM.leftAction, Turing.FinTM.leftCfg}
\difficulty{3}
Let $f : S \to S'$ be a state renaming, $a$ a one-step action of a machine
with $k$ work tapes over the Boolean alphabet, $c$ a configuration of such a
machine on input $x$, and fix contents and head positions for $l$ inactive
extra tapes.  Applying the left-block extension of $a$ along $f$ to the
left-block embedding of $c$ with the given inactive data yields the
left-block embedding, with the same inactive data, of the result of applying
$a$ to $c$.
\end{lemma}

\begin{lemma}[Action application commutes with the right embedding]
\lean{Turing.FinTM.rightCfg_apply}
\label{Turing.FinTM.rightCfg_apply}
\leanok
\uses{Turing.Action, Turing.Action.apply, Turing.Cfg, Turing.FinTM.rightAction, Turing.FinTM.rightCfg}
\difficulty{3}
Let $f : S \to S'$ be a state renaming, $a$ a one-step action of a machine
with $l$ work tapes over the Boolean alphabet, $c$ a configuration of such a
machine on input $x$, and fix contents and head positions for $k$ inactive
extra tapes.  Applying the right-block extension of $a$ along $f$ to the
right-block embedding of $c$ with the given inactive data yields the
right-block embedding, with the same inactive data, of the result of applying
$a$ to $c$.
\end{lemma}

\begin{lemma}[One-step simulation on the left block]
\lean{Turing.FinTM.leftCfg_step}
\label{Turing.FinTM.leftCfg_step}
\leanok
\uses{Turing.Action.apply, Turing.Cfg, Turing.Cfg.inputSymbol, Turing.Cfg.workTapeSymbols, Turing.FinTM.leftAction, Turing.FinTM.leftCfg, Turing.FinTM.leftCfg_apply, Turing.MultiTapeTM, Turing.MultiTapeTM.step}
\difficulty{4}
Let $M$ be a machine with $k$ work tapes and $M'$ a machine with $k + l$ work
tapes, both over the Boolean alphabet, and let $f$ be a map from the states
of $M$ to the states of $M'$ such that on every renamed state $f(q)$, for
every scanned input symbol and every tuple of scanned work symbols, the
transition of $M'$ is the left-block extension along $f$ of the transition of
$M$ from $q$ on that input symbol and the symbols scanned on the first $k$
tapes.  Then for every configuration $c$ of $M$ on input $x$ and all inactive
contents and head positions for the last $l$ tapes, one step of $M'$ from the
left-block embedding of $c$ equals the left-block embedding of one step of
$M$ from $c$; in particular a halted embedded configuration stays halted.
\end{lemma}

\begin{lemma}[One-step simulation on the right block]
\lean{Turing.FinTM.rightCfg_step}
\label{Turing.FinTM.rightCfg_step}
\leanok
\uses{Turing.Action.apply, Turing.Cfg, Turing.Cfg.inputSymbol, Turing.Cfg.workTapeSymbols, Turing.FinTM.rightAction, Turing.FinTM.rightCfg, Turing.FinTM.rightCfg_apply, Turing.MultiTapeTM, Turing.MultiTapeTM.step}
\difficulty{4}
Let $M$ be a machine with $l$ work tapes and $M'$ a machine with $k + l$ work
tapes, both over the Boolean alphabet, and let $f$ be a map from the states
of $M$ to the states of $M'$ such that on every renamed state $f(q)$, for
every scanned input symbol and every tuple of scanned work symbols, the
transition of $M'$ is the right-block extension along $f$ of the transition
of $M$ from $q$ on that input symbol and the symbols scanned on the last $l$
tapes.  Then for every configuration $c$ of $M$ on input $x$ and all inactive
contents and head positions for the first $k$ tapes --- which may hold
arbitrary data from an earlier phase --- one step of $M'$ from the
right-block embedding of $c$ equals the right-block embedding of one step of
$M$ from $c$.
\end{lemma}

\begin{lemma}[Lockstep run on the left block]
\lean{Turing.FinTM.leftCfg_run}
\label{Turing.FinTM.leftCfg_run}
\leanok
\uses{Turing.Cfg, Turing.FinTM.leftAction, Turing.FinTM.leftCfg, Turing.FinTM.leftCfg_step, Turing.MultiTapeTM, Turing.MultiTapeTM.runFrom, Turing.MultiTapeTM.runFrom_comm_of_step}
\difficulty{1}
Let $M$ be a machine with $k$ work tapes and $M'$ a machine with $k + l$ work
tapes, both over the Boolean alphabet, and let $f$ be a map from the states
of $M$ to the states of $M'$ such that on every renamed state, for all
scanned symbols, the transition of $M'$ is the left-block extension along $f$
of the corresponding transition of $M$ read off the first $k$ tapes.  Then
for every configuration $c$ of $M$ on input $x$, all inactive contents and
head positions for the last $l$ tapes, and every number of steps $t$, running
$M'$ for $t$ steps from the left-block embedding of $c$ equals the left-block
embedding of the run of $M$ for $t$ steps from $c$.
\end{lemma}

\begin{lemma}[Lockstep run on the right block]
\lean{Turing.FinTM.rightCfg_run}
\label{Turing.FinTM.rightCfg_run}
\leanok
\uses{Turing.Cfg, Turing.FinTM.branchTM, Turing.FinTM.rightAction, Turing.FinTM.rightCfg, Turing.FinTM.rightCfg_step, Turing.MultiTapeTM, Turing.MultiTapeTM.runFrom, Turing.MultiTapeTM.runFrom_comm_of_step}
\difficulty{1}
Let $M$ be a machine with $l$ work tapes and $M'$ a machine with $k + l$ work
tapes, both over the Boolean alphabet, and let $f$ be a map from the states
of $M$ to the states of $M'$ such that on every renamed state, for all
scanned symbols, the transition of $M'$ is the right-block extension along
$f$ of the corresponding transition of $M$ read off the last $l$ tapes.  Then
for every configuration $c$ of $M$ on input $x$, all inactive contents and
head positions for the first $k$ tapes, and every number of steps $t$,
running $M'$ for $t$ steps from the right-block embedding of $c$ equals the
right-block embedding of the run of $M$ for $t$ steps from $c$, with the
inactive left-block data preserved throughout.
\end{lemma}

\begin{definition}[Branch union of two machines]
\lean{Turing.FinTM.branchTM}
\label{Turing.FinTM.branchTM}
\leanok
\uses{Turing.FinTM, Turing.FinTM.leftAction, Turing.FinTM.rightAction}
Given two Turing machines $M_1$ and $M_2$ over the Boolean alphabet with
finitely many states, and a Boolean $b$, the machine whose work tapes are the
disjoint union of the tapes of $M_1$ and $M_2$, whose state type is the
disjoint union of their state types, whose transitions run $M_1$ on the left
tape block from left states and $M_2$ on the right tape block from right
states, and whose initial state is the initial state of $M_1$ when $b$ is
true and of $M_2$ when $b$ is false.  Only the initial state depends on $b$;
the transition table does not.
\end{definition}

\begin{lemma}[Branch union computes the selected branch]
\lean{Turing.FinTM.branchTM_computes}
\label{Turing.FinTM.branchTM_computes}
\leanok
\uses{Turing.FinTM, Turing.FinTM.ComputesInTime, Turing.FinTM.branchTM, Turing.FinTM.bufferTape, Turing.FinTM.computesInTime_iff, Turing.FinTM.leftCfg, Turing.FinTM.leftCfg_run, Turing.FinTM.rightCfg, Turing.FinTM.rightCfg_run, Turing.MultiTapeTM.initCfg}
\difficulty{4}
Let $M_1$ and $M_2$ be Turing machines over the Boolean alphabet with
finitely many states, let $b$ be a Boolean, let $x$ and $w$ be Boolean words,
and let $t$ be a natural number.  The branch union of $M_1$ and $M_2$
selected by $b$ halts on input $x$ within $t$ steps with output $w$ if and
only if the selected machine --- $M_1$ when $b$ is true, $M_2$ when $b$ is
false --- halts on input $x$ within $t$ steps with output $w$.
\end{lemma}

\begin{lemma}[Effect of a control action]
\lean{Turing.FinTM.controlAction_apply}
\label{Turing.FinTM.controlAction_apply}
\leanok
\uses{Turing.Action.apply, Turing.Cfg, Turing.FinTM.controlAction, Turing.moveInputPos}
\difficulty{2}
Let $c$ be a configuration of a machine with $k$ work tapes over the Boolean
alphabet on input $x$, let $m \in \{-1, 0, +1\}$ be a movement, and let $q$
be an optional state.  Applying the control action for $m$ and $q$ to $c$
yields the configuration that agrees with $c$ on all work-tape contents,
work-head positions, and output, whose state is $q$, and whose input head is
that of $c$ moved by $m$ with clamping at the two boundary blanks.
\end{lemma}

\begin{lemma}[Clamped left move subtracts one]
\lean{Turing.FinTM.moveInputPos_neg_val}
\label{Turing.FinTM.moveInputPos_neg_val}
\leanok
\uses{Turing.moveInputPos, Turing.moveInputPos_neg_of_ne_left}
\difficulty{2}
Let $p \in \{0,\dots,n+1\}$ be an input-head position for an input of length
$n$.  The clamped left move sends $p$ to the position whose value is $p - 1$
in truncated natural subtraction; in particular the left boundary position
$0$ is fixed.
\end{lemma}

\begin{lemma}[Leftward scan to the input start]
\lean{Turing.FinTM.rewind_scan}
\label{Turing.FinTM.rewind_scan}
\leanok
\uses{Turing.Cfg, Turing.Cfg.inputSymbol, Turing.FinTM.controlAction, Turing.FinTM.controlAction_apply, Turing.FinTM.moveInputPos_neg_val, Turing.MultiTapeTM, Turing.MultiTapeTM.runFrom, Turing.MultiTapeTM.runFrom_succ_eq_step, Turing.MultiTapeTM.step, Turing.inputSymbolInner, Turing.moveInputPos, Turing.moveInputPos_pos_of_ne_right}
\difficulty{5}
Let $M$ be a multi-tape Turing machine over the Boolean alphabet on input
$x$, let $s$ be a scanning state, and let $d$ be an optional dispatch state,
such that regardless of the work symbols read: in state $s$ scanning an input
symbol, $M$ moves the input head left and stays in $s$; in state $s$ scanning
a blank, $M$ moves the input head right and passes to $d$; and neither
transition touches a work tape or the output.  Then from any configuration
$c$ in state $s$ whose input-head position $p$ satisfies $p \le \abs{x}$,
running $M$ for exactly $p + 1$ steps yields the configuration that agrees
with $c$ in every component except that its state is $d$ and its input head
is at position $1$, the position scanning the first input symbol.
\end{lemma}

\begin{lemma}[Rewind from any input position]
\lean{Turing.FinTM.rewind_from_any}
\label{Turing.FinTM.rewind_from_any}
\leanok
\uses{Turing.Cfg, Turing.FinTM.controlAction, Turing.FinTM.controlAction_apply, Turing.FinTM.moveInputPos_neg_val, Turing.FinTM.rewind_scan, Turing.MultiTapeTM, Turing.MultiTapeTM.runFrom, Turing.MultiTapeTM.runFrom_add, Turing.MultiTapeTM.step, Turing.moveInputPos}
\difficulty{4}
Let $M$ be a multi-tape Turing machine over the Boolean alphabet on input
$x$, with a start state $r$, a scanning state $s$, and an optional dispatch
state $d$, such that regardless of the symbols read: from $r$ the machine
always moves the input head left and enters $s$; from $s$ it moves left and
stays in $s$ when an input symbol is scanned, and moves right and passes to
$d$ when a blank is scanned; and none of these transitions touches a work
tape or the output.  Then from any configuration $c$ in state $r$ there
exists a number of steps $t$ such that after $t$ steps the configuration
agrees with $c$ in every component except that its state is $d$ and its input
head is at position $1$.  This holds even when $x$ is empty or the input head
of $c$ starts at a boundary blank.
\end{lemma}

\begin{definition}[Three-block tape assembly]
\lean{Turing.FinTM.tapeBlocks}
\label{Turing.FinTM.tapeBlocks}
\leanok
Given a family of values indexed by $\{0,\dots,k-1\}$ (the left block), a
single value (the buffer), and a family of values indexed by
$\{0,\dots,l-1\}$ (the right block), the family indexed by the $k + (1 + l)$
tape indices that returns the left values on the first $k$ indices, the
buffer value on the next index, and the right values on the last $l$ indices.
\end{definition}

\begin{lemma}[Left projection of the tape assembly]
\lean{Turing.FinTM.tapeBlocks_left}
\label{Turing.FinTM.tapeBlocks_left}
\leanok
\uses{Turing.FinTM.tapeBlocks}
\difficulty{1}
Let $a$ be a family of values on $k$ indices, $b$ a single value, and $c$ a
family of values on $l$ indices.  Evaluating their three-block assembly at
the $i$-th index of the left block returns $a(i)$.
\end{lemma}

\begin{lemma}[Buffer projection of the tape assembly]
\lean{Turing.FinTM.tapeBlocks_buffer}
\label{Turing.FinTM.tapeBlocks_buffer}
\leanok
\uses{Turing.FinTM.tapeBlocks}
\difficulty{1}
Let $a$ be a family of values on $k$ indices, $b$ a single value, and $c$ a
family of values on $l$ indices.  Evaluating their three-block assembly at
the middle buffer index returns $b$.
\end{lemma}

\begin{lemma}[Right projection of the tape assembly]
\lean{Turing.FinTM.tapeBlocks_right}
\label{Turing.FinTM.tapeBlocks_right}
\leanok
\uses{Turing.FinTM.tapeBlocks}
\difficulty{1}
Let $a$ be a family of values on $k$ indices, $b$ a single value, and $c$ a
family of values on $l$ indices.  Evaluating their three-block assembly at
the $i$-th index of the right block returns $c(i)$.
\end{lemma}

\begin{definition}[Buffer tape storing a word]
\lean{Turing.FinTM.bufferTape}
\label{Turing.FinTM.bufferTape}
\leanok
For a Boolean word $w$, the work-tape content over the integer cells that
stores $w$ contiguously from cell zero: cell $z$ holds the $z$-th symbol of
$w$ (zero-indexed) when $0 \le z < \abs{w}$, and is blank at every other
integer cell.
\end{definition}

\begin{lemma}[Empty buffer is blank]
\lean{Turing.FinTM.bufferTape_nil}
\label{Turing.FinTM.bufferTape_nil}
\leanok
\uses{Turing.FinTM.bufferTape}
\difficulty{1}
The buffer tape storing the empty word is blank at every integer cell.
\end{lemma}

\begin{lemma}[Buffer reads at natural cells]
\lean{Turing.FinTM.bufferTape_nat}
\label{Turing.FinTM.bufferTape_nat}
\leanok
\uses{Turing.FinTM.bufferTape}
\difficulty{1}
Let $w$ be a Boolean word and $i$ a natural number.  Cell $i$ of the buffer
tape storing $w$ holds the optional $i$-th entry of $w$: the symbol $w_i$
when $i < \abs{w}$ and a blank otherwise.  In particular cell $\abs{w}$ is
the blank immediately to the right of the stored word.
\end{lemma}

\begin{lemma}[Left blank of a buffer]
\lean{Turing.FinTM.bufferTape_left}
\label{Turing.FinTM.bufferTape_left}
\leanok
\uses{Turing.FinTM.bufferTape}
\difficulty{1}
For every Boolean word $w$, cell $-1$ of the buffer tape storing $w$ is
blank; this includes the case of the empty word.
\end{lemma}

\begin{lemma}[Appending one bit to a buffer]
\lean{Turing.FinTM.bufferTape_append}
\label{Turing.FinTM.bufferTape_append}
\leanok
\uses{Turing.FinTM.bufferTape}
\difficulty{4}
Let $w$ be a Boolean word and $b$ a Boolean symbol.  The buffer tape storing
$w$ followed by the single symbol $b$ equals the buffer tape storing $w$
updated at cell $\abs{w}$ to hold $b$: exactly the cell that was the old
right blank changes, and every other integer cell is unaffected.
\end{lemma}

\begin{definition}[Boundary tag for a virtual head]
\lean{Turing.FinTM.VirtualTag}
\label{Turing.FinTM.VirtualTag}
\leanok
For a position $p \in \{0,\dots,n+1\}$ of a virtual input head over an input
of length $n$ and a Boolean tag $b$, the predicate asserting that $b$
correctly marks the boundaries: if $p = 0$ then $b$ is false, and if
$p = n + 1$ then $b$ is true.  At interior positions both tag values are
admissible.
\end{definition}

\begin{definition}[Clamped virtual head movement]
\lean{Turing.FinTM.virtualMove}
\label{Turing.FinTM.virtualMove}
\leanok
Given a Boolean boundary tag $b$, an optional scanned symbol, and a requested
movement $m \in \{-1, 0, +1\}$, the movement actually performed by a buffer
head simulating an input head: when the scanned symbol is blank and the
requested movement points outward --- left with tag false, marking the left
boundary, or right with tag true, marking the right boundary --- the head
stays put; in every other case it moves by $m$.
\end{definition}

\begin{definition}[Updated arrival tag]
\lean{Turing.FinTM.virtualNextTag}
\label{Turing.FinTM.virtualNextTag}
\leanok
Given a Boolean tag $b$ and a performed movement $m \in \{-1, 0, +1\}$, the
tag after the move: false after a left move, true after a right move, and $b$
itself after a stationary move, so that a suppressed outward move retains its
boundary tag and repeated outward attempts remain clamped.
\end{definition}

\begin{lemma}[Buffer reads agree with input reads]
\lean{Turing.FinTM.bufferTape_inputSymbol}
\label{Turing.FinTM.bufferTape_inputSymbol}
\leanok
\uses{Turing.Cfg, Turing.Cfg.inputSymbol, Turing.FinTM.bufferTape, Turing.FinTM.bufferTape_nat, Turing.FinTM.inputSymbol_at}
\difficulty{4}
Let $c$ be a configuration of a machine with $k$ work tapes over the Boolean
alphabet whose input is the word $w$, and let $p \in \{0,\dots,\abs{w}+1\}$
be its input-head position.  Cell $p - 1$ of the buffer tape storing $w$
holds exactly the symbol scanned by the input head of $c$: the corresponding
input symbol at interior positions, and a blank at the two boundary
positions.
\end{lemma}

\begin{lemma}[Virtual movement implements clamping]
\lean{Turing.FinTM.virtualMove_correct}
\label{Turing.FinTM.virtualMove_correct}
\leanok
\uses{Turing.Cfg, Turing.Cfg.inputSymbol, Turing.FinTM.VirtualTag, Turing.FinTM.bufferedCompTM, Turing.FinTM.virtualMove, Turing.FinTM.virtualNextTag, Turing.inputSymbolInner, Turing.moveInputPos, Turing.moveInputPos_neg_of_ne_left, Turing.moveInputPos_pos_of_ne_right}
\difficulty{5}
Let $c$ be a configuration of a machine over the Boolean alphabet on input
$w$, let $p \in \{0,\dots,\abs{w}+1\}$ be its input-head position, let $b$ be
a Boolean tag valid for $p$ (forced false at the left boundary $p = 0$ and
true at the right boundary $p = \abs{w} + 1$), and let $m \in \{-1, 0, +1\}$
be a requested movement.  Write $p'$ for the result of the native clamped
move of $p$ by $m$.  Then, as integers, $p - 1$ plus the virtual movement
determined by $b$, the scanned input symbol of $c$, and $m$ equals $p' - 1$;
and the updated tag --- false after a left move, true after a right move,
unchanged when stationary --- is again valid for $p'$.  This holds for every
word $w$, including the empty word.
\end{lemma}

\begin{definition}[Buffered sequential composition machine]
\lean{Turing.FinTM.bufferedCompTM}
\label{Turing.FinTM.bufferedCompTM}
\leanok
\uses{Turing.FinTM, Turing.FinTM.tapeBlocks, Turing.FinTM.virtualMove, Turing.FinTM.virtualNextTag}
Given machines $M_1$ and $M_2$ over the Boolean alphabet with finitely many
states, having $k_1$ and $k_2$ work tapes respectively, the machine with
$k_1 + (1 + k_2)$ work tapes that operates in phases.  In the first phase it
simulates $M_1$ on the left tape block, diverting every symbol that $M_1$
would output onto the middle buffer tape instead of the real output; the
simulated halt of $M_1$ leads to a live administrative rewind state, in which
the buffer head scans left until it reaches the blank preceding the stored
word.  It then moves right and dispatches into the second phase, a simulation
of $M_2$ on the right tape block in which the input reads of $M_2$ are served
from the buffer, with movements clamped at the boundary blanks of the stored
word by a Boolean arrival tag carried in the state.  Only the second phase
emits symbols to the real output or halts.
\end{definition}

\begin{definition}[First-phase embedded configuration]
\lean{Turing.FinTM.bufferedFirstCfg}
\label{Turing.FinTM.bufferedFirstCfg}
\leanok
\uses{Turing.Cfg, Turing.FinTM, Turing.FinTM.bufferTape, Turing.FinTM.bufferedCompTM, Turing.FinTM.tapeBlocks}
Given machines $M_1$ and $M_2$ over the Boolean alphabet with finitely many
states and a configuration $c$ of $M_1$ on input $x$, the configuration of
the buffered sequential composition of $M_1$ and $M_2$ on input $x$ whose
state is the first-phase image of the state of $c$ (a live administrative
state even when $c$ has halted), whose input head agrees with $c$, whose left
tape block carries the work tapes and heads of $c$, whose middle buffer tape
stores the output of $c$ contiguously with the buffer head on the blank
immediately to its right, whose right tape block is blank with heads at the
origin, and whose real output is empty.
\end{definition}

\begin{lemma}[Initialization of the buffered composition]
\lean{Turing.FinTM.bufferedFirstCfg_init}
\label{Turing.FinTM.bufferedFirstCfg_init}
\leanok
\uses{Turing.FinTM, Turing.FinTM.bufferedCompTM, Turing.FinTM.bufferedFirstCfg, Turing.FinTM.tapeBlocks, Turing.MultiTapeTM.initCfg}
\difficulty{3}
For all machines $M_1$ and $M_2$ over the Boolean alphabet with finitely many
states and every input $x$, the initial configuration of the buffered
sequential composition of $M_1$ and $M_2$ on $x$ equals the first-phase
embedding of the initial configuration of $M_1$ on $x$, in which the buffer
stores the empty word.
\end{lemma}

\begin{lemma}[First-phase single-step lockstep]
\lean{Turing.FinTM.bufferedFirstCfg_step}
\label{Turing.FinTM.bufferedFirstCfg_step}
\leanok
\uses{Turing.Action, Turing.Action.apply, Turing.Cfg, Turing.Cfg.inputSymbol, Turing.Cfg.workTapeSymbols, Turing.FinTM, Turing.FinTM.bufferTape_append, Turing.FinTM.bufferedCompTM, Turing.FinTM.bufferedFirstCfg, Turing.FinTM.tapeBlocks, Turing.MultiTapeTM.step}
\difficulty{5}
Let $M_1$ and $M_2$ be machines over the Boolean alphabet with finitely many
states, and let $c$ be a configuration of $M_1$ on input $x$ that has not
halted.  One step of the buffered sequential composition of $M_1$ and $M_2$
from the first-phase embedding of $c$ equals the first-phase embedding of one
step of $M_1$ from $c$.  In particular a symbol emitted by $M_1$ on that step
--- including on its halting transition --- lands on the buffer with the
buffer head advancing past it, the real output stays empty, and a simulated
halt of $M_1$ leaves the composite in a live administrative state.
\end{lemma}

\begin{lemma}[First-phase lockstep run]
\lean{Turing.FinTM.bufferedFirstCfg_run}
\label{Turing.FinTM.bufferedFirstCfg_run}
\leanok
\uses{Turing.Cfg, Turing.FinTM, Turing.FinTM.bufferedCompTM, Turing.FinTM.bufferedFirstCfg, Turing.FinTM.bufferedFirstCfg_step, Turing.MultiTapeTM.runFrom, Turing.MultiTapeTM.runFrom_succ_eq_step'}
\difficulty{3}
Let $M_1$ and $M_2$ be machines over the Boolean alphabet with finitely many
states, let $c$ be a configuration of $M_1$ on input $x$, and let $t$ be a
number of steps such that the run of $M_1$ from $c$ has not halted at any
time strictly before $t$.  Then running the buffered sequential composition
of $M_1$ and $M_2$ for $t$ steps from the first-phase embedding of $c$ yields
the first-phase embedding of the run of $M_1$ for $t$ steps from $c$.  The
time $t$ itself may be the first halting time of $M_1$.
\end{lemma}

\begin{definition}[Second-phase embedded configuration]
\lean{Turing.FinTM.bufferedSecondCfg}
\label{Turing.FinTM.bufferedSecondCfg}
\leanok
\uses{Turing.Cfg, Turing.FinTM, Turing.FinTM.bufferTape, Turing.FinTM.bufferedCompTM, Turing.FinTM.tapeBlocks}
Given machines $M_1$ and $M_2$ over the Boolean alphabet with finitely many
states, a configuration $c$ of $M_2$ on a virtual input $y$, a Boolean tag
$b$, a position $p \in \{0,\dots,\abs{x}+1\}$ of the physical input head over
the real input $x$, and arbitrary contents and head positions for the
left tape block of $M_1$, the configuration of the buffered sequential
composition of $M_1$ and $M_2$ on input $x$ whose state is the second-phase
image of the state of $c$ paired with $b$ (halted exactly when $c$ is
halted), whose physical input head is parked at $p$, whose left block carries
the given inactive data, whose middle buffer stores $y$ with the buffer head
at the virtual input-head position of $c$ minus one, whose right block
carries the work tapes and heads of $c$, and whose real output is the output
of $c$.
\end{definition}

\begin{lemma}[Second-phase single-step simulation]
\lean{Turing.FinTM.bufferedSecondCfg_step}
\label{Turing.FinTM.bufferedSecondCfg_step}
\leanok
\uses{Turing.Action.apply, Turing.Cfg, Turing.Cfg.inputSymbol, Turing.Cfg.workTapeSymbols, Turing.FinTM, Turing.FinTM.VirtualTag, Turing.FinTM.bufferTape_inputSymbol, Turing.FinTM.bufferedCompTM, Turing.FinTM.bufferedSecondCfg, Turing.FinTM.tapeBlocks_buffer, Turing.FinTM.virtualMove, Turing.FinTM.virtualMove_correct, Turing.FinTM.virtualNextTag, Turing.MultiTapeTM.step, Turing.MultiTapeTM.step_of_halt, Turing.moveInputPos_zero}
\difficulty{5}
Let $M_1$ and $M_2$ be machines over the Boolean alphabet with finitely many
states, let $c$ be a configuration of $M_2$ on a virtual input $y$, and let
$b$ be a Boolean tag valid for the virtual input-head position of $c$ (forced
false at the left boundary and true at the right boundary).  Fix a position
$p$ of the physical input head over the real input $x$ and arbitrary inactive
contents and head positions for the left tape block.  Then there exists a tag
$b'$, valid for the input-head position reached after one step of $M_2$ from
$c$, such that one step of the buffered sequential composition of $M_1$ and
$M_2$ from the second-phase embedding of $c$ with tag $b$ equals the
second-phase embedding of one step of $M_2$ from $c$ with tag $b'$, keeping
$p$ and the inactive data unchanged.  The halted case is included: a halted
$c$ stays halted on both sides.
\end{lemma}

\begin{lemma}[Second-phase lockstep run]
\lean{Turing.FinTM.bufferedSecondCfg_run}
\label{Turing.FinTM.bufferedSecondCfg_run}
\leanok
\uses{Turing.Cfg, Turing.FinTM, Turing.FinTM.VirtualTag, Turing.FinTM.bufferedCompTM, Turing.FinTM.bufferedSecondCfg, Turing.FinTM.bufferedSecondCfg_step, Turing.MultiTapeTM.runFrom, Turing.MultiTapeTM.runFrom_succ_eq_step'}
\difficulty{4}
Let $M_1$ and $M_2$ be machines over the Boolean alphabet with finitely many
states, let $c$ be a configuration of $M_2$ on a virtual input $y$, and let
$b$ be a Boolean tag valid for the virtual input-head position of $c$.  Fix a
position $p$ of the physical input head over the real input $x$, arbitrary
inactive contents and head positions for the left tape block, and a number of
steps $t$.  Then there exists a tag $b'$, valid for the input-head position
of the run of $M_2$ for $t$ steps from $c$, such that running the buffered
sequential composition of $M_1$ and $M_2$ for $t$ steps from the second-phase
embedding of $c$ with tag $b$ equals the second-phase embedding of that run
with tag $b'$, keeping $p$ and the inactive data unchanged.  Completed
outputs and absorbing halting are thereby preserved.
\end{lemma}

\begin{definition}[Rewind-scan configuration]
\lean{Turing.FinTM.bufferedScanCfg}
\label{Turing.FinTM.bufferedScanCfg}
\leanok
\uses{Turing.Cfg, Turing.FinTM, Turing.FinTM.bufferTape, Turing.FinTM.bufferedCompTM, Turing.FinTM.tapeBlocks}
Given machines $M_1$ and $M_2$ over the Boolean alphabet with finitely many
states, a Boolean word $y$ stored on the buffer, a position
$p \in \{0,\dots,\abs{x}+1\}$ of the physical input head over the real input
$x$, arbitrary contents and head positions for the left tape block, and a
natural number $j$, the configuration of the buffered sequential composition
of $M_1$ and $M_2$ on input $x$ in the live administrative scan state, with
the physical input head parked at $p$, the left block carrying the given
inactive data, the buffer storing $y$ with its head at cell $j - 1$, the
right block blank with heads at the origin, and empty real output.  The scan
state is live even when $j = 0$.
\end{definition}

\begin{lemma}[Scan reaches the initialized second phase]
\lean{Turing.FinTM.bufferedScanCfg_run}
\label{Turing.FinTM.bufferedScanCfg_run}
\leanok
\uses{Turing.Action.apply, Turing.Cfg.workTapeSymbols, Turing.FinTM, Turing.FinTM.bufferTape, Turing.FinTM.bufferTape_left, Turing.FinTM.bufferTape_nat, Turing.FinTM.bufferedCompTM, Turing.FinTM.bufferedScanCfg, Turing.FinTM.bufferedSecondCfg, Turing.FinTM.tapeBlocks_buffer, Turing.MultiTapeTM.initCfg, Turing.MultiTapeTM.runFrom, Turing.MultiTapeTM.runFrom_succ_eq_step, Turing.MultiTapeTM.runFrom_succ_eq_step', Turing.MultiTapeTM.runFrom_zero, Turing.MultiTapeTM.step, Turing.moveInputPos_zero}
\difficulty{5}
Let $M_1$ and $M_2$ be machines over the Boolean alphabet with finitely many
states, let $y$ be a Boolean word stored on the buffer, let $p$ be a position
of the physical input head over the real input $x$, and fix arbitrary
inactive contents and head positions for the left tape block.  For every
$j \le \abs{y}$, running the buffered sequential composition of $M_1$ and
$M_2$ for exactly $j + 1$ steps from the rewind-scan configuration with
buffer head at cell $j - 1$ yields the second-phase embedding of the initial
configuration of $M_2$ on virtual input $y$, with arrival tag true, physical
position $p$, and the same inactive data.  This includes the case of an
empty word $y$.
\end{lemma}

\begin{lemma}[Exact cost of rewind and dispatch]
\lean{Turing.FinTM.bufferedFirstCfg_rewind}
\label{Turing.FinTM.bufferedFirstCfg_rewind}
\leanok
\uses{Turing.Action.apply, Turing.Cfg, Turing.FinTM, Turing.FinTM.bufferedCompTM, Turing.FinTM.bufferedFirstCfg, Turing.FinTM.bufferedScanCfg, Turing.FinTM.bufferedScanCfg_run, Turing.FinTM.bufferedSecondCfg, Turing.MultiTapeTM.initCfg, Turing.MultiTapeTM.runFrom, Turing.MultiTapeTM.runFrom_succ_eq_step, Turing.MultiTapeTM.step, Turing.moveInputPos_zero}
\difficulty{4}
Let $M_1$ and $M_2$ be machines over the Boolean alphabet with finitely many
states, and let $c$ be a halted configuration of $M_1$ on input $x$ with
output $w$.  Running the buffered sequential composition of $M_1$ and $M_2$
for exactly $\abs{w} + 2$ steps from the first-phase embedding of $c$ yields
the second-phase embedding of the initial configuration of $M_2$ on virtual
input $w$, with arrival tag true, the physical input head parked where $c$
left it, and the work tapes and head positions of $c$ retained as inactive
left-block data.  The real output remains empty throughout.
\end{lemma}

\begin{lemma}[Reaching the second phase from the start]
\lean{Turing.FinTM.bufferedComp_start}
\label{Turing.FinTM.bufferedComp_start}
\leanok
\uses{Turing.FinTM, Turing.FinTM.ComputesInTime, Turing.FinTM.ComputesInTime.output_unique, Turing.FinTM.bufferedCompTM, Turing.FinTM.bufferedFirstCfg_init, Turing.FinTM.bufferedFirstCfg_rewind, Turing.FinTM.bufferedFirstCfg_run, Turing.FinTM.bufferedSecondCfg, Turing.FinTM.computesInTime_iff, Turing.MultiTapeTM.initCfg, Turing.MultiTapeTM.runFrom, Turing.MultiTapeTM.runFrom_add}
\difficulty{5}
Let $M_1$ and $M_2$ be machines over the Boolean alphabet with finitely many
states, and suppose that $M_1$, on input $x$, halts within $t_1$ steps with
output $y$.  Then there exist a step count $a \le t_1 + \abs{y} + 2$, a
position $p$ of the physical input head over $x$, and contents and head
positions for the left tape block, such that running the buffered sequential
composition of $M_1$ and $M_2$ for $a$ steps from its initial configuration
on $x$ yields the second-phase embedding of the initial configuration of
$M_2$ on virtual input $y$, with arrival tag true, physical position $p$, and
that inactive left-block data.
\end{lemma}
